% Copy the contents into an existing chapter file; keep its source label.
% Chapter numbers are local to the volume, assigned by its main .tex file.
\chapter{A chapter title that makes a promise}
\label{ch:replace-with-existing-source-label}

\ChapterLead{One or two sentences establishing the chapter's central question.}

% Begin with a concrete situation and establish the problem in connected prose.
% Introduce the necessary context here; do not assume another book was read.
% Optional: a compact route through a long lesson, after its hook or lead.
% \begin{LessonBrief}
% \BriefPoint{The question}{The consequential question this lesson answers.}
% \BriefPoint{The insight}{The mechanism that changes the reader's understanding.}
% \BriefPoint{The decision}{The choice that understanding makes possible.}
% \end{LessonBrief}

\section{The first turn in the argument}
\label{sec:replace-with-existing-section-label}

% Develop the claim, mechanism, example, evidence, and limit in paragraphs.
% A section should answer a question, not list everything related to a noun.
% Use subsections only when the argument needs an additional level.

\section{What changes when the conditions change}

% A counterexample or boundary makes the argument useful outside one case.
% Use a table or diagram only when it clarifies a relationship described here.

\section{Making the decision}

% Apply the argument to the continuing case and show the actual choice.
% Close with the decision now possible, then transition to the next chapter.

% Available asides (use sparingly):
% \begin{NoteBox}[A useful qualification] ... \end{NoteBox}
% \begin{TipBox}[Try this] ... \end{TipBox}
% \begin{DefinitionBox}{The term} ... \end{DefinitionBox}
% \begin{WarningBox}{A tempting shortcut} ... \end{WarningBox}
% \begin{KeyIdea}One short central claim.\end{KeyIdea}
% Emphasis is upright semibold. Do not introduce EB Garamond italics.

% Captionable diagram example. Explain it in the prose before using it.
% \begin{figure}[htbp]
% \centering
% \begin{tikzpicture}[node distance=12mm,
%   every node/.style={draw,align=center,font=\sffamily\small}]
% \node (a) {Starting condition};
% \node[right=of a] (b) {Mechanism};
% \node[right=of b] (c) {Changed outcome};
% \draw[->] (a)--(b); \draw[->,dashed] (b)--(c);
% \end{tikzpicture}
% \caption{Explain the relationship and what the dashed link means.}
% \label{fig:unique-label}
% \FigureSource{Actual source and context, or an explicit illustrative-data statement.}
% \end{figure}
% Use Figure~\ref{fig:unique-label} in the prose.
% The design specimen includes reusable bar/pie/radar/waterfall examples.
% PGFPlots is loaded: use \begin{axis}[BIchart,...] inside tikzpicture.

% Tables use tabularx/booktabs, caption, and label inside a table environment.
% Use \sffamily\small for dense data; keep ordinary text at the reading size.
% Equations use equation/aligned or \BookEquation{...}; label meaningful results.

% A short worked example stays together and receives a chapter-local number.
% \begin{WorkedExample}{An explanatory title}
% \label{ex:unique-label}
% \WorkedStep{1}{Set the boundary.}{State the inputs, units, and assumptions.}
% \WorkedStep{2}{Do the work.}{Show the reasoning or calculation.}
% \WorkedStep{3}{Interpret it.}{Explain the result and its limits.}
% \end{WorkedExample}
% Refer to Worked Example~\ref{ex:unique-label} in the prose.
% For a genuinely long example, use \begin{WorkedExample}[breakable]{Title}.

% Place an evidence note beside the claim it qualifies.
% \begin{EvidenceNote}{The study or dataset}
% \EvidenceLine{Source}{Actual citation, population, sample size, and period.}
% \EvidenceLine{Finding}{The supported result, with units and uncertainty.}
% \EvidenceLine{Limit}{What the design cannot establish or generalise to.}
% \end{EvidenceNote}
% For invented data, state explicitly that no empirical study was conducted.

% A short decision checklist fits before the cases, not after every section.
% \begin{DecisionChecklist}{A specific decision}
% \item A concrete fact, calculation, or condition the reader can verify.
% \item An assumption to test before committing resources.
% \item Evidence that would change the decision.
% \end{DecisionChecklist}

\LessonTakeaway{The decision this lesson now makes possible.}

\CaseStudies{Introduce the shared decision. These are fictional teaching cases;
each founder operates under different constraints.}

\CaseStudy{1}{TECHNOLOGY}{A title that names the decision}
  {Founder name\quad /\quad Business context\quad /\quad Stage}
% Establish a concrete situation in prose.
\CaseSection{The decision}
% Explain the alternatives and the binding constraint.
\CasePrompt{Ask the reader to choose and state their reasoning.}
\CaseSection{A defensible move}
% Apply the chapter's mechanism, including evidence and tradeoffs.
\CaseSection{What would change the answer?}
% Give a boundary or changed condition, not a universal prescription.

\CaseStudy{2}{FASHION}{A title that names the decision}
  {Founder name\quad /\quad Business context\quad /\quad Stage}
% Follow the same reasoning moves with genuinely different constraints.

\CaseStudy{3}{EDUCATION}{A title that names the decision}
  {Founder name\quad /\quad Business context\quad /\quad Stage}
% Show how the mechanism changes when the promised result is capability.

% Aim for 250-300 words per short case; allow longer cases to flow.
% An optional captioned comparison table can connect the three decisions.

% Optional closing recap after the cases. Avoid repeating a lesson brief verbatim.
% \begin{LessonReview}
% \item One transferable mechanism, stated in a complete sentence.
% \item A boundary or failure condition.
% \item The decision now possible.
% \end{LessonReview}
% \TransferQuestion{A question applying the lesson to the reader's own business.}

\NotesSection

% Add real research notes and sources. Distinguish observed facts from
% hypothetical figures and the author's judgment. Never invent citations.
